%% =====================================================================
%%  B5-T-05-03 -- what B5 contributes to "The Black Swan"
%%  -------------------------------------------------------------------
%%  LAYER A: APPEND-ONLY. Written once, then left alone. Material found
%%  only here sits in the `unique` slot and is protected by rule R10.
%% =====================================================================
%% @kind: source
%% @id: B5-T-05-03
%% @book: B5
%% @topic: T-05-03
%% @locator: ch. 10, pp. 281--312
%% @status: distilled
%% @unique: yes
%% @integrated: yes
%% @updated: 2026-09-19
%% @allow-rewrite: B5 ch. 8-10 ingest 2026-09-19 -- committed scaffold replaced by the written entry (planned -> stable lifecycle)

\dossier{B5}{T-05-03}{\bookshort{B5}}{ch. 10, pp. 281--312}

\begin{distilled}
  \item The chapter imports the black swan idea: what made swans surprising was not that they were unknowable but that observers had sampled only white ones --- negotiation runs on the same unexamined induction.
  \item Unknown unknowns are named as the class whose unearthing has game-changing effects: breakthroughs that shift the game inalterably come from identifying and using them.
  \item The cost of assumption is stated from the hostage canon: over-reliance on known knowns shackles a negotiator to assumptions that prevent seeing and hearing what the situation actually presents.
  \item The similarity principle: research confirms we trust people more when they read as similar or familiar; belonging is a primal instinct, and triggering the sense of \emph{we see the world the same way} immediately gains ground.
  \item The cases turn on facts nobody thought to ask about --- a hostage-taker's specific grievance, a hidden stake --- which once spoken re-priced the entire negotiation.
  \item The discipline prescribed is heavy, patient listening at close range: face time, curiosity about the counterpart's world, and letting them correct your picture of it.
\end{distilled}

\begin{unique}
  \item The three-part epistemology (known knowns, unknown knowns, unknown unknowns) as a working checklist for preparation: write what you assume, and treat each assumption as a question you have not asked yet.
  \item The similarity-contact pair as the disclosure mechanism: swans surface only when the counterpart feels in-group and stays in rich channel --- which makes rapport a data-collection instrument, not a nicety.
  \item The re-pricing doctrine: a black swan is not added to the deal, it changes what the deal is.
\end{unique}
